\documentclass[10pt,a4paper]{article}
\usepackage[margin=2.2cm]{geometry}
\usepackage{fontspec}
\setmonofont{DejaVu Sans Mono}[Scale=0.84]
\newfontfamily\unifont{DejaVu Serif}[Scale=0.9]
\usepackage{amsmath,amssymb,mathtools}
\setcounter{MaxMatrixCols}{20}
\usepackage{longtable,booktabs,array}
\usepackage{fancyvrb}
\usepackage[dvipsnames]{xcolor}
\usepackage[colorlinks=true,linkcolor=NavyBlue,urlcolor=NavyBlue]{hyperref}
\usepackage{microtype}
\usepackage{needspace}
\usepackage{tocloft}\setlength{\cftbeforesecskip}{0.45em}
\setlength{\parindent}{0pt}\setlength{\parskip}{0.55em}
\setlength{\emergencystretch}{8em}
\renewcommand{\arraystretch}{1.15}
\allowdisplaybreaks
\begin{document}
\title{dirac16complex: symbols and their definitions\\[0.4em]\large A dictionary of every symbol used for the Dirac equation and the Kohn–Sham theories of the 16-component quantized fermion field dirac16complex on 7+1 spacetime — sections 7–10 of the notebook 7+1.nb}
\author{Repository: github.com/once-ere/7plus1 — prepared for Nash\\ Notebook: 7+1.nb (Wolfram Language 15.0.1)\\ Document generated by tools/make\_docs.py on 2026-10-27}
\date{}
\maketitle
\begin{abstract}
For each symbol this dictionary gives the mathematical notation, the name used in the Mathematica notebook 7+1.nb (where there is one), the defining equation, and its meaning. The symbols are grouped by topic in the order in which they are introduced; longer defining equations follow each table as displayed formulas. Section 10 lists symbols that look alike, and section 11 lists every Mathematica name defined in sections 7–10 of 7+1.nb with its Wolfram Language definition. Every equation needed to define a symbol is written in this document.
\end{abstract}
\tableofcontents
\Needspace{12\baselineskip}\section*{1. Conventions}\addcontentsline{toc}{section}{1. Conventions}
\begin{itemize}
\item Units: $\hbar=c=1$.
\item Index ranges: $A,B,K,L\in\lbrace0,1,\dots,7\rbrace$ (spacetime), $i,j\in\lbrace1,\dots,7\rbrace$ (space), $a,b,c,d\in\lbrace1,\dots,16\rbrace$ (spinor components), $s,r\in\lbrace1,\dots,8\rbrace$ (spin labels), $k,l$ (orbitals or Grassmann generators). Repeated indices are summed unless stated otherwise.
\item Indices are raised and lowered with the metric $\eta$ of section 2: $a_A=\eta_{AA}a^A$ (no sum), so $a_0=a^0$ and $a_i=-a^i$.
\item Matrices: $1_n$ is the $n\times n$ identity, $M^{\mathsf T}$ the transpose, $\overline M$ or $M^*$ the complex conjugate, $M^\dagger=\overline M^{\mathsf T}$ the Hermitian conjugate, $\mathrm{tr}$ the trace, $\delta_{ij}$ the Kronecker delta, $\lbrace X,Y\rbrace=XY+YX$ the anticommutator and $[X,Y]=XY-YX$ the commutator.
\item Block notation: $\left(\begin{smallmatrix}a&b\\c&d\end{smallmatrix}\right)$ is the 16×16 matrix with the four 8×8 blocks $a,b,c,d$; a 16-component column is written $(\psi_\uparrow;\psi_\downarrow)$, upper 8 components first.
\item Typefaces: Kohn–Sham potentials are calligraphic ($\mathcal S,\mathcal P,\mathcal V_A,\mathcal A_A$), external sources are sans-serif ($\mathsf s,\mathsf p,\mathsf v_A,\mathsf a_A$). Mathematica names are set in \texttt{typewriter}.
\item Mathematica lists are 1-based: $\eta_{AB}$ is \texttt{η17[[A+1, B+1]]}, $\psi_a$ is \texttt{dirac16complex[[a]]}. Greek letters in names are Mathematica named characters: \texttt{\textbackslash{}[Psi]} is \texorpdfstring{\ensuremath{\psi}}{ψ}, \texttt{\textbackslash{}[Phi]} is \texorpdfstring{\ensuremath{\phi}}{ϕ}, \texttt{\textbackslash{}[CapitalGamma]} is \texorpdfstring{\ensuremath{\Gamma}}{Γ}, and so on.
\item Wirtinger convention: a complex variable $z$ and its conjugate $z^*$ are treated as independent variables when differentiating, $\partial z/\partial z=1$, $\partial z^*/\partial z=0$.
\end{itemize}
\Needspace{12\baselineskip}\section*{2. Spacetime and momenta}\addcontentsline{toc}{section}{2. Spacetime and momenta}
{\small
\begin{longtable}{>{\raggedright\arraybackslash}p{\dimexpr 0.1500\linewidth-2\tabcolsep\relax}>{\raggedright\arraybackslash}p{\dimexpr 0.1900\linewidth-2\tabcolsep\relax}>{\raggedright\arraybackslash}p{\dimexpr 0.3600\linewidth-2\tabcolsep\relax}>{\raggedright\arraybackslash}p{\dimexpr 0.3000\linewidth-2\tabcolsep\relax}}
\toprule
\textbf{symbol} & \textbf{Mathematica} & \textbf{definition} & \textbf{meaning} \\
\midrule\endfirsthead
\toprule
\textbf{symbol} & \textbf{Mathematica} & \textbf{definition} & \textbf{meaning} \\
\midrule\endhead
\endfoot
\bottomrule\endlastfoot
$X$; $x_0,\dots,x_7$ & \texttt{X = \{x0, …, x7\}} & $X=(x_0,x_1,\dots,x_7)$, $x^A=x_A$ as symbols & $x_0$ ordinary time (direction of the octonion unit $e_0=1$); $x_1,x_2,x_3$ ordinary 3-space; $x_4,\dots,x_7$ auxiliary space-like \\
$x$, $\vec x$ & — & $x=(x_0,\vec x)$, $\vec x=(x_1,\dots,x_7)$ & spacetime point; space part \\
$\eta_{AB}$, $\eta^{AB}$ & \texttt{η17}, \texttt{η17inv} & $\eta=\mathrm{diag}(+1,-1,-1,-1,-1,-1,-1,-1)$, $\eta^{-1}=\eta$ & metric of 7+1 spacetime \\
$a\cdot b$ & — & $\eta_{AB}a^Ab^B=a^0b^0-\sum_ia^ib^i$ & Minkowski product \\
$\partial_A$ & \texttt{D[f, X[[A+1]]]} & $\partial_A=\partial/\partial x_A$, $\partial_Ax^B=\delta_A^B$ & partial derivative \\
$\Box$ & \texttt{waveOp17} & $\eta^{AB}\partial_A\partial_B=\partial_0^2-\partial_1^2-\dots-\partial_7^2$ & wave (d’Alembert) operator \\
$\vec p$, $p^i$ & \texttt{pv7} & $\vec p=(p^1,\dots,p^7)$; \texttt{pv7 = \{p1, …, p7\}} & momentum; real symbols in the symbolic checks \\
$\vec p\cdot\vec x$, $\vec p^{\,2}$ & — & $\sum_ip^ix_i$, $\sum_i(p^i)^2$ & Euclidean products of space vectors \\
$E$ & \texttt{Ep} & $E=\sqrt{m^2+\vec p^{\,2}}$ & energy of a free particle of momentum $\vec p$ \\
$p$ & — & $p=(E,\vec p)$, $p\cdot p=m^2$ & 8-momentum on the mass shell \\
$e_s$ & \texttt{UnitVector[8, s]} & $(e_s)_r=\delta_{rs}$ & unit column of length 8 \\
$\delta^7(\vec x-\vec y)$ & — & $\int\delta^7(\vec x-\vec y)f(\vec y)\,d^7y=f(\vec x)$ & Dirac delta in 7 space dimensions \\*
\end{longtable}}
The Minkowski square of a point is the real part of the square of the octonion $x=x_0e_0+\sum_ix_ie_i$:
\begin{equation*}
x\cdot x=\eta_{AB}x^Ax^B=x_0^2-x_1^2-\dots-x_7^2=\mathrm{Re}(x\,x).
\end{equation*}
\Needspace{12\baselineskip}\section*{3. The real 8×8 matrices}\addcontentsline{toc}{section}{3. The real 8×8 matrices}
{\small
\begin{longtable}{>{\raggedright\arraybackslash}p{\dimexpr 0.1500\linewidth-2\tabcolsep\relax}>{\raggedright\arraybackslash}p{\dimexpr 0.1900\linewidth-2\tabcolsep\relax}>{\raggedright\arraybackslash}p{\dimexpr 0.3600\linewidth-2\tabcolsep\relax}>{\raggedright\arraybackslash}p{\dimexpr 0.3000\linewidth-2\tabcolsep\relax}}
\toprule
\textbf{symbol} & \textbf{Mathematica} & \textbf{definition} & \textbf{meaning} \\
\midrule\endfirsthead
\toprule
\textbf{symbol} & \textbf{Mathematica} & \textbf{definition} & \textbf{meaning} \\
\midrule\endhead
\endfoot
\bottomrule\endlastfoot
$\tau^{O}_i$ & \texttt{τO[i]} & the seven real 8×8 matrices displayed below; $\tau^{O}_0=1_8$ & left multiplication by the imaginary octonion unit $e_i$, written in a spinor basis \\
$\tau_A$ & \texttt{τ[A]} & $\tau_0=1_8$; $\tau_i=\tau^{O}_i$ ($i\le3$); $\tau_4,\dots,\tau_7$ displayed below & the real split-octonion matrices of the notebook 4+4.nb \\
$\sigma$ & \texttt{σ} & $\sigma=\tau_1\tau_2\tau_3=\left(\begin{smallmatrix}0&1_4\\1_4&0\end{smallmatrix}\right)$ & the 8×8 matrix with $\tau^{O}_k=\tau_k\sigma$ for $k=4,\dots,7$ \\
$\bar\tau_A$ & \texttt{τbar[A]} & $\bar\tau_0=1_8$, $\bar\tau_A=\sigma\,\tau_A^{\mathsf T}\sigma$ ($A=1,\dots,7$) & conjugate split-octonion matrices of 4+4.nb \\
$\eta^{(4,4)}$ & \texttt{η4488} & $\mathrm{diag}(+1,+1,+1,+1,-1,-1,-1,-1)$ & metric of the notebook 4+4.nb \\
$T_A$ & \texttt{T16[A]} & $T_A=\left(\begin{smallmatrix}0&\bar\tau_A\\\tau_A&0\end{smallmatrix}\right)$, $\lbrace T_A,T_B\rbrace=2\eta^{(4,4)}_{AB}1_{16}$ & the 16×16 Dirac matrices of 4+4.nb \\
$\sigma_i$ & \texttt{σW[i]} & $\sigma_i=i\,\tau^{O}_i$ & Hermitian; the Pauli matrices of 7 space dimensions \\
$\tau^{O}\cdot\vec p$ & — & $\sum_ip^i\tau^{O}_i$ & real, antisymmetric \\
$\sigma\cdot\vec p$ & \texttt{σWp[p]} & $\sum_ip^i\sigma_i$ & Hermitian \\
$\sigma^A$, $\bar\sigma^A$ & \texttt{σUp}, \texttt{σbarUp} & $\sigma^A=(1_8,\sigma_1,\dots,\sigma_7)$, $\bar\sigma^A=(1_8,-\sigma_1,\dots,-\sigma_7)$ & chiral 8×8 matrices (lists of 8 matrices; \texttt{σUp[[A+1]]} is $\sigma^A$) \\*
\end{longtable}}
Defining relations ($i,j=1,\dots,7$):
\begin{align*}
(\tau^{O}_i)^{\mathsf T}&=-\tau^{O}_i,\qquad \tau^{O}_i\tau^{O}_j+\tau^{O}_j\tau^{O}_i=-2\,\delta_{ij}\,1_8,\qquad \tau^{O}_1\tau^{O}_2\cdots\tau^{O}_7=1_8, \\
\sigma_i^\dagger&=\sigma_i,\qquad \sigma_i\sigma_j+\sigma_j\sigma_i=2\,\delta_{ij}\,1_8,\qquad (\tau^{O}\cdot\vec p)^2=-\vec p^{\,2}\,1_8,\qquad (\sigma\cdot\vec p)^2=\vec p^{\,2}\,1_8, \\
\sigma^2&=1_8,\qquad \sigma\tau_k=\tau_k\sigma\ (k=1,2,3),\qquad \sigma\tau_k=-\tau_k\sigma,\ \ \tau_k^2=+1_8\ (k=4,\dots,7).
\end{align*}
The seven matrices $\tau^{O}_i$:
\begin{equation*}
\tau^{O}_1=\begin{pmatrix}0 & 0 & 0 & 0 & 0 & 0 & 0 & 1 \\ 0 & 0 & 0 & 0 & 0 & 0 & 1 & 0 \\ 0 & 0 & 0 & 0 & 0 & -1 & 0 & 0 \\ 0 & 0 & 0 & 0 & -1 & 0 & 0 & 0 \\ 0 & 0 & 0 & 1 & 0 & 0 & 0 & 0 \\ 0 & 0 & 1 & 0 & 0 & 0 & 0 & 0 \\ 0 & -1 & 0 & 0 & 0 & 0 & 0 & 0 \\ -1 & 0 & 0 & 0 & 0 & 0 & 0 & 0\end{pmatrix},\quad \tau^{O}_2=\begin{pmatrix}0 & 0 & 0 & 0 & 0 & 0 & -1 & 0 \\ 0 & 0 & 0 & 0 & 0 & 0 & 0 & 1 \\ 0 & 0 & 0 & 0 & 1 & 0 & 0 & 0 \\ 0 & 0 & 0 & 0 & 0 & -1 & 0 & 0 \\ 0 & 0 & -1 & 0 & 0 & 0 & 0 & 0 \\ 0 & 0 & 0 & 1 & 0 & 0 & 0 & 0 \\ 1 & 0 & 0 & 0 & 0 & 0 & 0 & 0 \\ 0 & -1 & 0 & 0 & 0 & 0 & 0 & 0\end{pmatrix},
\end{equation*}
\begin{equation*}
\tau^{O}_3=\begin{pmatrix}0 & 0 & 0 & 0 & 0 & 1 & 0 & 0 \\ 0 & 0 & 0 & 0 & -1 & 0 & 0 & 0 \\ 0 & 0 & 0 & 0 & 0 & 0 & 0 & 1 \\ 0 & 0 & 0 & 0 & 0 & 0 & -1 & 0 \\ 0 & 1 & 0 & 0 & 0 & 0 & 0 & 0 \\ -1 & 0 & 0 & 0 & 0 & 0 & 0 & 0 \\ 0 & 0 & 0 & 1 & 0 & 0 & 0 & 0 \\ 0 & 0 & -1 & 0 & 0 & 0 & 0 & 0\end{pmatrix},\quad \tau^{O}_4=\begin{pmatrix}0 & 1 & 0 & 0 & 0 & 0 & 0 & 0 \\ -1 & 0 & 0 & 0 & 0 & 0 & 0 & 0 \\ 0 & 0 & 0 & -1 & 0 & 0 & 0 & 0 \\ 0 & 0 & 1 & 0 & 0 & 0 & 0 & 0 \\ 0 & 0 & 0 & 0 & 0 & -1 & 0 & 0 \\ 0 & 0 & 0 & 0 & 1 & 0 & 0 & 0 \\ 0 & 0 & 0 & 0 & 0 & 0 & 0 & 1 \\ 0 & 0 & 0 & 0 & 0 & 0 & -1 & 0\end{pmatrix},
\end{equation*}
\begin{equation*}
\tau^{O}_5=\begin{pmatrix}0 & 0 & -1 & 0 & 0 & 0 & 0 & 0 \\ 0 & 0 & 0 & -1 & 0 & 0 & 0 & 0 \\ 1 & 0 & 0 & 0 & 0 & 0 & 0 & 0 \\ 0 & 1 & 0 & 0 & 0 & 0 & 0 & 0 \\ 0 & 0 & 0 & 0 & 0 & 0 & 1 & 0 \\ 0 & 0 & 0 & 0 & 0 & 0 & 0 & 1 \\ 0 & 0 & 0 & 0 & -1 & 0 & 0 & 0 \\ 0 & 0 & 0 & 0 & 0 & -1 & 0 & 0\end{pmatrix},\quad \tau^{O}_6=\begin{pmatrix}0 & 0 & 0 & -1 & 0 & 0 & 0 & 0 \\ 0 & 0 & 1 & 0 & 0 & 0 & 0 & 0 \\ 0 & -1 & 0 & 0 & 0 & 0 & 0 & 0 \\ 1 & 0 & 0 & 0 & 0 & 0 & 0 & 0 \\ 0 & 0 & 0 & 0 & 0 & 0 & 0 & 1 \\ 0 & 0 & 0 & 0 & 0 & 0 & -1 & 0 \\ 0 & 0 & 0 & 0 & 0 & 1 & 0 & 0 \\ 0 & 0 & 0 & 0 & -1 & 0 & 0 & 0\end{pmatrix},
\end{equation*}
\begin{equation*}
\tau^{O}_7=\begin{pmatrix}0 & 0 & 0 & 0 & -1 & 0 & 0 & 0 \\ 0 & 0 & 0 & 0 & 0 & -1 & 0 & 0 \\ 0 & 0 & 0 & 0 & 0 & 0 & -1 & 0 \\ 0 & 0 & 0 & 0 & 0 & 0 & 0 & -1 \\ 1 & 0 & 0 & 0 & 0 & 0 & 0 & 0 \\ 0 & 1 & 0 & 0 & 0 & 0 & 0 & 0 \\ 0 & 0 & 1 & 0 & 0 & 0 & 0 & 0 \\ 0 & 0 & 0 & 1 & 0 & 0 & 0 & 0\end{pmatrix}.
\end{equation*}
The split-octonion matrices $\tau_4,\dots,\tau_7$ of 4+4.nb (they square to $+1_8$; $\tau^{O}_k=\tau_k\,\sigma$):
\begin{equation*}
\tau_4=\begin{pmatrix}0 & 0 & 0 & 0 & 0 & 1 & 0 & 0 \\ 0 & 0 & 0 & 0 & -1 & 0 & 0 & 0 \\ 0 & 0 & 0 & 0 & 0 & 0 & 0 & -1 \\ 0 & 0 & 0 & 0 & 0 & 0 & 1 & 0 \\ 0 & -1 & 0 & 0 & 0 & 0 & 0 & 0 \\ 1 & 0 & 0 & 0 & 0 & 0 & 0 & 0 \\ 0 & 0 & 0 & 1 & 0 & 0 & 0 & 0 \\ 0 & 0 & -1 & 0 & 0 & 0 & 0 & 0\end{pmatrix},\quad \tau_5=\begin{pmatrix}0 & 0 & 0 & 0 & 0 & 0 & -1 & 0 \\ 0 & 0 & 0 & 0 & 0 & 0 & 0 & -1 \\ 0 & 0 & 0 & 0 & 1 & 0 & 0 & 0 \\ 0 & 0 & 0 & 0 & 0 & 1 & 0 & 0 \\ 0 & 0 & 1 & 0 & 0 & 0 & 0 & 0 \\ 0 & 0 & 0 & 1 & 0 & 0 & 0 & 0 \\ -1 & 0 & 0 & 0 & 0 & 0 & 0 & 0 \\ 0 & -1 & 0 & 0 & 0 & 0 & 0 & 0\end{pmatrix},
\end{equation*}
\begin{equation*}
\tau_6=\begin{pmatrix}0 & 0 & 0 & 0 & 0 & 0 & 0 & -1 \\ 0 & 0 & 0 & 0 & 0 & 0 & 1 & 0 \\ 0 & 0 & 0 & 0 & 0 & -1 & 0 & 0 \\ 0 & 0 & 0 & 0 & 1 & 0 & 0 & 0 \\ 0 & 0 & 0 & 1 & 0 & 0 & 0 & 0 \\ 0 & 0 & -1 & 0 & 0 & 0 & 0 & 0 \\ 0 & 1 & 0 & 0 & 0 & 0 & 0 & 0 \\ -1 & 0 & 0 & 0 & 0 & 0 & 0 & 0\end{pmatrix},\quad \tau_7=\begin{pmatrix}-1 & 0 & 0 & 0 & 0 & 0 & 0 & 0 \\ 0 & -1 & 0 & 0 & 0 & 0 & 0 & 0 \\ 0 & 0 & -1 & 0 & 0 & 0 & 0 & 0 \\ 0 & 0 & 0 & -1 & 0 & 0 & 0 & 0 \\ 0 & 0 & 0 & 0 & 1 & 0 & 0 & 0 \\ 0 & 0 & 0 & 0 & 0 & 1 & 0 & 0 \\ 0 & 0 & 0 & 0 & 0 & 0 & 1 & 0 \\ 0 & 0 & 0 & 0 & 0 & 0 & 0 & 1\end{pmatrix} .
\end{equation*}
\Needspace{12\baselineskip}\section*{4. Dirac matrices, chirality and the chiral projections}\addcontentsline{toc}{section}{4. Dirac matrices, chirality and the chiral projections}
{\small
\begin{longtable}{>{\raggedright\arraybackslash}p{\dimexpr 0.1500\linewidth-2\tabcolsep\relax}>{\raggedright\arraybackslash}p{\dimexpr 0.1900\linewidth-2\tabcolsep\relax}>{\raggedright\arraybackslash}p{\dimexpr 0.3600\linewidth-2\tabcolsep\relax}>{\raggedright\arraybackslash}p{\dimexpr 0.3000\linewidth-2\tabcolsep\relax}}
\toprule
\textbf{symbol} & \textbf{Mathematica} & \textbf{definition} & \textbf{meaning} \\
\midrule\endfirsthead
\toprule
\textbf{symbol} & \textbf{Mathematica} & \textbf{definition} & \textbf{meaning} \\
\midrule\endhead
\endfoot
\bottomrule\endlastfoot
$\Gamma_A$ & \texttt{Γ17[A]} & $\Gamma_0=\left(\begin{smallmatrix}0&1_8\\1_8&0\end{smallmatrix}\right)$, $\Gamma_i=\left(\begin{smallmatrix}\tau^{O}_i&0\\0&-\tau^{O}_i\end{smallmatrix}\right)$ & the eight REAL 16×16 Dirac matrices; every entry is $0$ or $\pm1$ \\
$\Gamma^A$ & \texttt{Γ17up[A]} & $\Gamma^A=\eta^{AA}\Gamma_A$: $\Gamma^0=\Gamma_0$, $\Gamma^i=-\Gamma_i$ & Dirac matrix with an upper index \\
$\Gamma_K$ & \texttt{products256} & $\Gamma_K=\Gamma_{A_1}\Gamma_{A_2}\cdots\Gamma_{A_k}$ for $K=\lbrace A_1<\dots<A_k\rbrace$; $\Gamma_\emptyset=1_{16}$ & the 256 basis elements of $Cl(1,7)\cong\mathbb R(16)$ \\
$\beta$ & \texttt{βD} & $\beta=\Gamma^0=\left(\begin{smallmatrix}0&1_8\\1_8&0\end{smallmatrix}\right)$ & Dirac $\beta$; real symmetric \\
$\alpha_i$ & \texttt{αD[i]} & $\alpha_i=\Gamma^0\Gamma^i=\left(\begin{smallmatrix}0&\tau^{O}_i\\-\tau^{O}_i&0\end{smallmatrix}\right)$ & Dirac $\alpha$; real symmetric \\
$\omega$ & \texttt{ω17}, \texttt{Γ17[8]} & $\omega=\Gamma_0\Gamma_1\cdots\Gamma_7=\left(\begin{smallmatrix}0&-1_8\\1_8&0\end{smallmatrix}\right)$ & volume element; real, $\omega^2=-1_{16}$ \\
$\Sigma_{AB}$ & \texttt{Σ17[A, B]} & $\Sigma_{AB}=\tfrac14(\Gamma_A\Gamma_B-\Gamma_B\Gamma_A)$ & generators of Spin(1,7); real \\
$\chi$ & \texttt{χ17} & $\chi=i\,\omega=\left(\begin{smallmatrix}0&-i1_8\\i1_8&0\end{smallmatrix}\right)$ & chirality matrix; Hermitian, complex \\
$P_L$, $P_R$ & \texttt{PL17}, \texttt{PR17} & $P_L=\tfrac12(1_{16}-\chi)$, $P_R=\tfrac12(1_{16}+\chi)$ & chiral projections; complex, rank 8 each \\
$E_L$, $E_R$ & \texttt{EL}, \texttt{ER} & $E_L=\tfrac1{\sqrt2}\left(\begin{smallmatrix}1_8\\-i1_8\end{smallmatrix}\right)$, $E_R=\tfrac1{\sqrt2}\left(\begin{smallmatrix}-i1_8\\1_8\end{smallmatrix}\right)$ & 16×8 isometries onto the eigenspaces $\chi=-1$ and $\chi=+1$ \\
$W$ & \texttt{Wch} & $W=(E_L,E_R)$ & unitary 16×16 change to the chiral basis \\
$\eta'$ & — & $\eta'=-\eta=\mathrm{diag}(-1,+1,\dots,+1)$ & the opposite sign convention; generates $Cl(7,1)\cong\mathbb H(8)$, which has no real 16×16 realization \\
$R_A$, $T$, $Y$ & — & hypothetical real $R_A$ with $R_AR_B+R_BR_A=2\eta'_{AB}1_{16}$; $R_A=T(i\Gamma_A)T^{-1}$; $Y=T^{-1}\overline T$ & symbols of the proof that $\eta'$ has no real 16×16 realization ($Y$ would satisfy $Y\overline Y=1_{16}$ and $Y=c\,\omega$, impossible since $\omega^2=-1_{16}$) \\
$Z$ & — & real 16×16 matrix with $Z\Gamma_A=-\Gamma_AZ$ ($A=0,\dots,7$) & the solutions are $Z=c\,\omega$, $c\in\mathbb R$, and $Z^2=-c^2 1_{16}$: there is no real chirality matrix \\*
\end{longtable}}
Defining relations ($A,B=0,\dots,7$; $i,j=1,\dots,7$):
\begin{align*}
\Gamma_A\Gamma_B+\Gamma_B\Gamma_A&=2\,\eta_{AB}\,1_{16},\qquad \Gamma_0^2=1_{16},\quad \Gamma_i^2=-1_{16},\qquad \Gamma^A\Gamma_A=8\cdot1_{16}, \\
(\Gamma^A)^\dagger&=(\Gamma^A)^{\mathsf T}=\Gamma^0\Gamma^A\Gamma^0, \\
\alpha_i\alpha_j+\alpha_j\alpha_i&=2\,\delta_{ij}\,1_{16},\qquad \alpha_i\beta+\beta\alpha_i=0,\qquad \beta^2=1_{16}, \\
\omega\,\Gamma_A&=-\Gamma_A\,\omega,\qquad \chi^\dagger=\chi,\qquad \chi^2=1_{16},\qquad \chi\,\Gamma_A=-\Gamma_A\,\chi,\qquad \chi\,\Sigma_{AB}=\Sigma_{AB}\,\chi, \\
P_L+P_R&=1_{16},\qquad P_L^2=P_L,\quad P_R^2=P_R,\quad P_LP_R=0,\qquad \Gamma_AP_L=P_R\Gamma_A,\qquad P_L^\dagger\Gamma^0=\Gamma^0P_R, \\
E_L^\dagger E_L&=E_R^\dagger E_R=1_8,\quad E_L^\dagger E_R=0,\quad E_LE_L^\dagger=P_L,\quad E_RE_R^\dagger=P_R, \\
\chi E_L&=-E_L,\quad \chi E_R=E_R,\quad W^\dagger W=1_{16}, \\
\mathrm{tr}\big(\Gamma_K^{-1}\Gamma_L\big)&=16\,\delta_{KL},\qquad \Gamma_K^{-1}=\Gamma_K^{\mathsf T}, \\
\textstyle\sum_K(\Gamma_K)_{ab}(\Gamma_K^{-1})_{cd}&=16\,\delta_{ad}\,\delta_{cb}\qquad(\text{Fierz completeness}).
\end{align*}
Blocks of the Dirac matrices between the chiral subspaces:
\begin{align*}
E_L^\dagger\beta E_R&=E_R^\dagger\beta E_L=1_8,\qquad E_L^\dagger\beta E_L=E_R^\dagger\beta E_R=0, \\
E_L^\dagger\alpha_iE_L&=-\sigma_i,\qquad E_R^\dagger\alpha_iE_R=\sigma_i,\qquad E_L^\dagger\alpha_iE_R=0 .
\end{align*}
\Needspace{12\baselineskip}\section*{5. The field dirac16complex and its free dynamics}\addcontentsline{toc}{section}{5. The field dirac16complex and its free dynamics}
{\small
\begin{longtable}{>{\raggedright\arraybackslash}p{\dimexpr 0.1500\linewidth-2\tabcolsep\relax}>{\raggedright\arraybackslash}p{\dimexpr 0.1900\linewidth-2\tabcolsep\relax}>{\raggedright\arraybackslash}p{\dimexpr 0.3600\linewidth-2\tabcolsep\relax}>{\raggedright\arraybackslash}p{\dimexpr 0.3000\linewidth-2\tabcolsep\relax}}
\toprule
\textbf{symbol} & \textbf{Mathematica} & \textbf{definition} & \textbf{meaning} \\
\midrule\endfirsthead
\toprule
\textbf{symbol} & \textbf{Mathematica} & \textbf{definition} & \textbf{meaning} \\
\midrule\endhead
\endfoot
\bottomrule\endlastfoot
$\psi$, $\psi_a$ & \texttt{dirac16complex}, \texttt{ψ16[a]} & $\psi=(\psi_1,\dots,\psi_{16})^{\mathsf T}$; \texttt{dirac16complex = Table[ψ16[a] @@ X, \{a, 1, 16\}]} & the quantized 16-component fermion field: complex Grassmann-odd operators \\
$\psi_\uparrow$, $\psi_\downarrow$ & parts \texttt{[[1;;8]]} and \texttt{[[9;;16]]} of \texttt{dirac16complex} & $\psi=(\psi_\uparrow;\psi_\downarrow)$ & upper and lower 8 components \\
$\bar\psi$ & — & $\bar\psi=\psi^\dagger\Gamma^0$ & Dirac adjoint \\
$m$ & \texttt{mD}; \texttt{1} in the examples & $m>0$ & mass of the free field \\
$\Gamma^A\partial_A\psi$ & \texttt{DiracOp17[ψ]} & $\sum_A\Gamma^A\partial_A\psi$ & Dirac operator \\
$(i\Gamma^A\partial_A-m)\psi$ & \texttt{DiracEq16[ψ, m]} & $i\Gamma^A\partial_A\psi-m\psi$ & left side of the Dirac equation \\
$h$ & \texttt{hD[ψ, m]} & $h=-i\alpha_i\partial_i+\beta m$ & Dirac Hamiltonian: $i\partial_0\psi=h\psi$ \\
$h(\vec p)$ & \texttt{hp[p, m]} & $h(\vec p)=p^i\alpha_i+m\beta$ & Dirac Hamiltonian of a plane wave; real symmetric \\
$M(\vec p)$ & \texttt{Mp}, \texttt{M0}, \texttt{M5} & $M(\vec p)=m1_8+p^i\tau^{O}_i$ & real 8×8 block of $h(\vec p)$ \\
$u_s(\vec p)$ & \texttt{uS[s, M, E]} & $u_s=\frac1{\sqrt E}(E\,e_s;\ M^{\mathsf T}e_s)$ & real positive-energy spinors \\
$v_s(\vec p)$ & \texttt{vS[s, M, E]} & $v_s=\frac1{\sqrt E}(E\,e_s;\ -M e_s)$ & real negative-energy spinors \\
$\not p$ & \texttt{slash17[E, p]} & $\not p=\Gamma_Ap^A=E\,\Gamma_0+p^i\Gamma_i$ & real 16×16 matrix \\
$U(\vec p)$ & \texttt{Umat} & columns $u_s(\vec p)/\sqrt{2E}$ and $v_s(-\vec p)/\sqrt{2E}$, $s=1,\dots,8$ & real orthogonal matrix diagonalizing $h(\vec p)$ \\
$j^A$ & \texttt{jG} (Grassmann check) & $j^A=\bar\psi\Gamma^A\psi$ & vector current; $j^0=\psi^\dagger\psi$, $j^i=\psi^\dagger\alpha_i\psi$ \\
$\mathcal L_0$ & — & $\mathcal L_0=\bar\psi(i\Gamma^A\partial_A-m)\psi$ & free Lagrangian density \\
$V$ & \texttt{Vbox} & — & volume of the periodic box of the mode expansion \\
$b_s(\vec p)$, $d_s(\vec p)$ & \texttt{modeB}, \texttt{modeD} (coefficients) & operators of the mode expansion below & annihilate a particle, an antiparticle \\
$H$, $Q$ & — & normal-ordered Hamiltonian and charge, below & energy and charge operators \\
$\psi^{c}$ & — & $\psi^{c}=\omega\,\psi^*$ & charge conjugate field \\
$\beta_{\rm b}$ & \texttt{β} (inside \texttt{Sβ}, \texttt{Λboost}) & real number & rapidity of a boost along $x_1$ \\
$\mathcal B(\beta_{\rm b})$ & \texttt{Sβ}, \texttt{Sβinv} & $\mathcal B=e^{\beta_{\rm b}\Sigma_{01}}=\cosh\frac{\beta_{\rm b}}2\,1_{16}+\sinh\frac{\beta_{\rm b}}2\,\Gamma_0\Gamma_1$ & spinor boost; \texttt{Sβinv} is $\mathcal B^{-1}$ \\
$L$ & \texttt{Λboost} & $L_{BK}=\frac1{16}\mathrm{tr}\big(\eta_{BB}\Gamma_B\mathcal B\Gamma_K\mathcal B^{-1}\big)$ & Lorentz matrix of the boost \\
$p'$ & \texttt{pPrime} & $p'=Lp$ & boosted 8-momentum \\
$\phi_L$, $\phi_R$ & \texttt{ϕL}, \texttt{ϕR} & $\phi_L=E_L^\dagger\psi$, $\phi_R=E_R^\dagger\psi$ & chiral (Weyl) 8-component fields \\
$\kappa$ (field) & — & $\kappa=\phi_L^\dagger\phi_R$ & complex scalar bilinear of the field \\*
\end{longtable}}
Equations that define and relate these symbols:
\begin{align*}
&\big(i\,\Gamma^A\partial_A-m\big)\psi=0,\qquad \big(i\Gamma^A\partial_A+m\big)\big(i\Gamma^A\partial_A-m\big)=-\big(\Box+m^2\big)1_{16},\qquad i\,\partial_0\psi=h\,\psi, \\
&h(\vec p)=\begin{pmatrix}0&M\\M^{\mathsf T}&0\end{pmatrix},\qquad MM^{\mathsf T}=M^{\mathsf T}M=E^2\,1_8,\qquad M(-\vec p)=M(\vec p)^{\mathsf T},\qquad h(\vec p)^2=E^2\,1_{16}, \\
&h(\vec p)\,u_s=E\,u_s,\qquad h(-\vec p)\,v_s=-E\,v_s,\qquad u_s^{\mathsf T}u_r=v_s^{\mathsf T}v_r=2E\,\delta_{sr},\qquad (\not p-m)u_s=0,\qquad (\not p+m)v_s=0, \\
&\textstyle\sum_su_s\bar u_s=\not p+m,\qquad \sum_sv_s\bar v_s=\not p-m,\qquad U^{\mathsf T}U=1_{16},\qquad U^{\mathsf T}h(\vec p)\,U=\mathrm{diag}(E\,1_8,-E\,1_8), \\
&\partial_Aj^A=0,\qquad \psi^{c}\text{ solves the Dirac equation when }\psi\text{ does},\qquad (\psi^c)^c=-\psi,\qquad L^{\mathsf T}\eta L=\eta .
\end{align*}
Mode expansion in the box of volume $V$, anticommutators, Hamiltonian and charge:
\begin{align*}
\psi(x_0,\vec x)&=\sum_{\vec p}\sum_{s=1}^{8}\frac{1}{\sqrt{2EV}}\Big[b_s(\vec p)\,u_s(\vec p)\,e^{-i(Ex_0-\vec p\cdot\vec x)}+d_s^\dagger(\vec p)\,v_s(\vec p)\,e^{+i(Ex_0-\vec p\cdot\vec x)}\Big], \\
\big\lbrace b_s(\vec p),b_r^\dagger(\vec q)\big\rbrace&=\big\lbrace d_s(\vec p),d_r^\dagger(\vec q)\big\rbrace=\delta_{sr}\,\delta_{\vec p\vec q},\qquad \big\lbrace\psi_a(x_0,\vec x),\psi_b^\dagger(x_0,\vec y)\big\rbrace=\delta_{ab}\,\delta^7(\vec x-\vec y), \\
H&=\int:\psi^\dagger h\,\psi:\,d^7x=\sum_{\vec p,s}E\,\big(b_s^\dagger b_s+d_s^\dagger d_s\big),\qquad Q=\int:\psi^\dagger\psi:\,d^7x=\sum_{\vec p,s}\big(b_s^\dagger b_s-d_s^\dagger d_s\big).
\end{align*}
Chiral form of the field equations and the bilinears ($\psi=E_L\phi_L+E_R\phi_R$):
\begin{align*}
\big(i\,\partial_0+\tau^{O}_i\,\partial_i\big)\,\phi_L&=m\,\phi_R,\qquad \big(i\,\partial_0-\tau^{O}_i\,\partial_i\big)\,\phi_R=m\,\phi_L,\qquad i\,\bar\sigma^A\partial_A\phi_L=m\,\phi_R,\quad i\,\sigma^A\partial_A\phi_R=m\,\phi_L, \\
\phi_L&=\frac{\psi_\uparrow+i\,\psi_\downarrow}{\sqrt2},\qquad \phi_R=\frac{i\,\psi_\uparrow+\psi_\downarrow}{\sqrt2},\qquad \bar\psi\psi=\kappa+\kappa^*,\qquad \bar\psi\,i\chi\,\psi=i(\kappa-\kappa^*) .
\end{align*}
\Needspace{12\baselineskip}\section*{6. Grassmann variables}\addcontentsline{toc}{section}{6. Grassmann variables}
{\small
\begin{longtable}{>{\raggedright\arraybackslash}p{\dimexpr 0.1500\linewidth-2\tabcolsep\relax}>{\raggedright\arraybackslash}p{\dimexpr 0.1900\linewidth-2\tabcolsep\relax}>{\raggedright\arraybackslash}p{\dimexpr 0.3600\linewidth-2\tabcolsep\relax}>{\raggedright\arraybackslash}p{\dimexpr 0.3000\linewidth-2\tabcolsep\relax}}
\toprule
\textbf{symbol} & \textbf{Mathematica} & \textbf{definition} & \textbf{meaning} \\
\midrule\endfirsthead
\toprule
\textbf{symbol} & \textbf{Mathematica} & \textbf{definition} & \textbf{meaning} \\
\midrule\endhead
\endfoot
\bottomrule\endlastfoot
$\theta_k$ & \texttt{θg[k]} $=$ \texttt{grass[\{k\}]} & $\theta_k\theta_l=-\theta_l\theta_k$, hence $\theta_k^2=0$ & Grassmann generators \\
$\theta_{k_1}\cdots\theta_{k_r}$ & \texttt{grass[\{k1, …, kr\}]} & ordered monomial, $k_1<\dots<k_r$; \texttt{grass[\{\}]} $=1$ & basis of the Grassmann algebra \\
product & \texttt{gMul[x, y]} & bilinear extension of the product of monomials below & Grassmann product \\
$\partial^{L}/\partial\theta_k$ & \texttt{gD[k, x]} & left derivative, below & moves $\theta_k$ to the front and removes it \\
$\bar\theta_k$ & \texttt{grass[\{3 + k\}]} & $\bar\theta_k=\theta_{3+k}$ & conjugate generators of the current-conservation check \\
$U_1,U_2,U_3$ & \texttt{Ugr} & three plane-wave solutions, below & ingredients of a Grassmann-valued solution \\
$\psi_G$, $\bar\psi_G$ & \texttt{ψG}, \texttt{ψbarG} & $\psi_G=\sum_k\theta_kU_k$, $\bar\psi_G=\sum_k\bar\theta_kU_k^\dagger\Gamma^0$ & Grassmann-valued solution and its adjoint \\
$\psi_a$, $\bar\psi_a$ (32 generators) & \texttt{ψP}, \texttt{ψbarP} & $\psi_a=\theta_a$, $\bar\psi_a=\theta_{16+a}$ & generators of the nonlinear-equation check \\
$\bar\psi M\psi$ & \texttt{gBil[M]} & $\sum_{a,b}\bar\psi_aM_{ab}\psi_b$ & Grassmann bilinear \\
$\mathcal L_{\rm int}$ & \texttt{LintG} & $-\frac{g_S}2(\bar\psi\psi)^2-\frac{g_P}2(\bar\psi i\chi\psi)^2-\frac{g_V}2(\bar\psi\Gamma^A\psi)(\bar\psi\Gamma_A\psi)-\frac{g_A}2(\bar\psi\Gamma^A\chi\psi)(\bar\psi\Gamma_A\chi\psi)$ & interaction Lagrangian (section 7) \\*
\end{longtable}}
Product and left derivative of monomials ($K$, $K'$ ordered index sets; $\pi$ the permutation that sorts the concatenation; the hat means omission):
\begin{align*}
\theta_K\,\theta_{K^{\prime}}&=\begin{cases}0,&K\cap K^{\prime}\neq\emptyset,\\ \mathrm{sign}(\pi)\,\theta_{\mathrm{sort}(K\cup K^{\prime})},&\text{otherwise},\end{cases} \\
\frac{\partial^{L}}{\partial\theta_k}\big(\theta_{k_1}\cdots\theta_{k_r}\big)&=(-1)^{j-1}\,\theta_{k_1}\cdots\widehat{\theta_{k_j}}\cdots\theta_{k_r}\quad\text{if }k=k_j\text{ (and }0\text{ if }k\notin\lbrace k_1,\dots,k_r\rbrace).
\end{align*}
The three plane waves ($m=1$):
\begin{align*}
U_1&=u_1(\vec p_1)\,e^{-i(5x_0-\vec p_1\cdot\vec x)},\qquad U_2=v_3(\vec p_3)\,e^{+i(3x_0-\vec p_3\cdot\vec x)},\qquad U_3=u_5(\vec 0)\,e^{-ix_0}, \\
\vec p_1&=(2,2,4,0,0,0,0),\qquad \vec p_3=(0,0,0,2,2,0,0).
\end{align*}
\Needspace{12\baselineskip}\section*{7. The self-coupled field and the 16-component Kohn–Sham theory}\addcontentsline{toc}{section}{7. The self-coupled field and the 16-component Kohn–Sham theory}
{\small
\begin{longtable}{>{\raggedright\arraybackslash}p{\dimexpr 0.1500\linewidth-2\tabcolsep\relax}>{\raggedright\arraybackslash}p{\dimexpr 0.1900\linewidth-2\tabcolsep\relax}>{\raggedright\arraybackslash}p{\dimexpr 0.3600\linewidth-2\tabcolsep\relax}>{\raggedright\arraybackslash}p{\dimexpr 0.3000\linewidth-2\tabcolsep\relax}}
\toprule
\textbf{symbol} & \textbf{Mathematica} & \textbf{definition} & \textbf{meaning} \\
\midrule\endfirsthead
\toprule
\textbf{symbol} & \textbf{Mathematica} & \textbf{definition} & \textbf{meaning} \\
\midrule\endhead
\endfoot
\bottomrule\endlastfoot
$g_S,g_P,g_V,g_A$ & \texttt{gS}, \texttt{gP}, \texttt{gV}, \texttt{gA} & real constants of mass dimension $-6$ & scalar, pseudoscalar, vector, axial-vector couplings; $g<0$ attractive \\
$\rho_S$ & \texttt{dens16[Γ17[0]]} & $\rho_S=\bar\psi\psi$ & scalar density \\
$\rho_P$ & \texttt{dens16[M]}, $M=\Gamma^0i\chi$ & $\rho_P=\bar\psi\,i\chi\,\psi$ & pseudoscalar density \\
$j^A$, $j_A$ & \texttt{dens16[M]}, $M=\Gamma^0\Gamma^A$ & $j^A=\bar\psi\Gamma^A\psi$, $j_A=\eta_{AA}j^A$ & vector current \\
$j_5^A$, $j_{5A}$ & \texttt{dens16[M]}, $M=\Gamma^0\Gamma^A\chi$ & $j_5^A=\bar\psi\Gamma^A\chi\psi$, $j_{5A}=\eta_{AA}j_5^A$ & axial-vector current \\
$\mathsf s,\mathsf p,\mathsf v_A,\mathsf a_A$ & — & real functions of $x$ & external sources coupled to $\rho_S$, $\rho_P$, $j^A$, $j_5^A$ \\
$n$ & — & $n=(\rho_S,\rho_P,j^A,j_5^A)$ & the 18 densities ($1+1+8+8$) \\
$w$ & — & $w=(\mathsf s,\mathsf p,\mathsf v_A,\mathsf a_A)$ & the 18 sources \\
$w\cdot n$ & — & $\mathsf s\rho_S+\mathsf p\rho_P+\mathsf v_Aj^A+\mathsf a_Aj_5^A$ & source term \\
$(g,O,O')$ & \texttt{vPairs} & $(g_S,1,1)$, $(g_P,i\chi,i\chi)$, $(g_V,\Gamma^A,\Gamma_A)$, $(g_A,\Gamma^A\chi,\Gamma_A\chi)$, $A=0,\dots,7$ & the 18 vertex triples ($1+1+8+8$): coupling and the two matrices of a contact term $-\frac g2(\bar\psi O\psi)(\bar\psi O'\psi)$ \\
$\mathcal L$ & — & below & Lagrangian density of the self-coupled field \\
$\mathcal H$, $H[w]$, $H_0$, $H_{\rm int}$ & — & below & Hamiltonian density and Hamiltonians \\
$N$ & — & $Q=N$ & number of particles (charge sector) \\
$F[n]$ & — & $F[n]=\min_{|\Psi\rangle\to n}\langle\Psi|H_0+H_{\rm int}|\Psi\rangle$ & universal (Levy–Lieb) functional \\
$E[w]$ & — & $E[w]=\min_n\big(F[n]+\int w\cdot n\,d^7x\big)$ & ground-state energy \\
$\varphi_k$, $\varepsilon_k$ & \texttt{orbV}, \texttt{orbcV} & $h_{\rm KS}\varphi_k=\varepsilon_k\varphi_k$, $\int\varphi_k^\dagger\varphi_l=\delta_{kl}$ & Kohn–Sham orbitals and eigenvalues; \texttt{orbV}, \texttt{orbcV} are $\varphi$, $\varphi^*$ as independent variables \\
$\mathcal S,\mathcal P$ & \texttt{Sks}, \texttt{Pks} & below & scalar and pseudoscalar Kohn–Sham potentials \\
$\mathcal V_A,\mathcal A_A$ & \texttt{VKS}, \texttt{AKS} & below; LOWER index & vector and axial Kohn–Sham potentials; \texttt{VKS[[A+1]]} is $\mathcal V_A$ \\
$h_{\rm KS}$ & \texttt{hKS[p]} & below & Kohn–Sham Hamiltonian; Hermitian for real potentials \\
$T_s$ & — & $T_s=\sum_k\int\varphi_k^\dagger(-i\alpha_i\partial_i+\beta m)\varphi_k\,d^7x$ & kinetic plus rest energy of the orbitals \\
$E_H$, $e_H$ & \texttt{EH16}, \texttt{EHfun} & below & Hartree energy and its density \\
$v_H$ & \texttt{vH16} & below & Hartree potential matrix \\
$E_{xc}$ & — & $E_{xc}=F-T_s-E_H$ & exchange-correlation energy \\
$E_x$, $E_c$ & \texttt{Exfun[R]} & $E_x$ below; $E_c=E_{xc}-E_x$ & exact exchange energy; correlation energy \\
$R$ & \texttt{Rorb}, \texttt{Rr}, \texttt{Runi} & $R=\sum_k\varphi_k\bar\varphi_k$ (16×16) & Kohn–Sham density matrix at a point \\
$\varrho$ & — & $\varrho_{xy}=\langle c_y^\dagger c_x\rangle=\sum_k\varphi_{k,x}\varphi_{k,y}^*$, $R=\varrho\,\Gamma^0$ & one-body density matrix \\
$\Sigma_x$ & \texttt{Sigmax[R]} & below & exchange self-energy \\
$\Sigma_H$ & — & $\Sigma_H=\mathcal S+i\chi\mathcal P+\Gamma^A(\mathcal V_A+\chi\mathcal A_A)$ & Hartree self-energy \\
$c_a$, $c_a^\dagger$ & \texttt{cdc[a, b]} $=c_a^\dagger c_b$ & $\lbrace c_a,c_b^\dagger\rbrace=\delta_{ab}$ & mode operators of the 16 components at one point \\
$B(M)$ & \texttt{Bop[M]} & $B(M)=\sum_{a,b}M_{ab}\,c_a^\dagger c_b$ & one-body operator, $\psi^\dagger M\psi$ \\
$|\Phi\rangle$ & \texttt{Φ3} & Slater determinant of 3 orthonormal orbitals & state of the Wick-theorem check (\texttt{orbs3}; Fock space of dimension $\binom{16}{3}=560$) \\*
\end{longtable}}
The Lagrangian and Hamiltonian densities, and the Hamiltonian with static sources:
\begin{align*}
\mathcal L&=\bar\psi\big(i\Gamma^A\partial_A-m\big)\psi-\bar\psi\big(\mathsf s+i\chi\,\mathsf p+\Gamma^A(\mathsf v_A+\chi\,\mathsf a_A)\big)\psi \\
&\quad-\frac{g_S}{2}(\bar\psi\psi)^2-\frac{g_P}{2}(\bar\psi\,i\chi\psi)^2-\frac{g_V}{2}(\bar\psi\Gamma^A\psi)(\bar\psi\Gamma_A\psi)-\frac{g_A}{2}(\bar\psi\Gamma^A\chi\psi)(\bar\psi\Gamma_A\chi\psi), \\
\mathcal H&=\,:\psi^\dagger\big(-i\alpha_i\partial_i+\beta m\big)\psi:+:\bar\psi\big(\mathsf s+i\chi\,\mathsf p+\Gamma^A(\mathsf v_A+\chi\,\mathsf a_A)\big)\psi: \\
&\quad+\frac{g_S}{2}:\rho_S^2:+\frac{g_P}{2}:\rho_P^2:+\frac{g_V}{2}:j^Aj_A:+\frac{g_A}{2}:j_5^Aj_{5A}:, \\
H[w]&=H_0+H_{\rm int}+\int w\cdot n\,d^7x,\qquad H_0=\int:\psi^\dagger\big(-i\alpha_i\partial_i+\beta m\big)\psi:\,d^7x .
\end{align*}
The Kohn–Sham Hamiltonian, potentials, densities and the Hartree quantities:
\begin{align*}
h_{\rm KS}&=\alpha_i\big(-i\partial_i+\mathcal V_i\big)+\beta\big(m+\mathcal S\big)+\mathcal V_0+\beta\,i\chi\,\mathcal P+\chi\,\mathcal A_0+\alpha_i\chi\,\mathcal A_i, \\
\mathcal S&=\mathsf s+g_S\,\rho_S+\frac{\delta E_{xc}}{\delta\rho_S},\qquad \mathcal P=\mathsf p+g_P\,\rho_P+\frac{\delta E_{xc}}{\delta\rho_P}, \\
\mathcal V_A&=\mathsf v_A+g_V\,j_A+\frac{\delta E_{xc}}{\delta j^A},\qquad \mathcal A_A=\mathsf a_A+g_A\,j_{5A}+\frac{\delta E_{xc}}{\delta j_5^A}, \\
\rho_S&=\sum_k\bar\varphi_k\varphi_k,\qquad \rho_P=\sum_k\bar\varphi_k\,i\chi\,\varphi_k,\qquad j^A=\sum_k\bar\varphi_k\Gamma^A\varphi_k,\qquad j_5^A=\sum_k\bar\varphi_k\Gamma^A\chi\,\varphi_k, \\
E_H&=\int e_H\,d^7x,\qquad e_H=\frac{g_S}{2}\rho_S^2+\frac{g_P}{2}\rho_P^2+\frac{g_V}{2}j^Aj_A+\frac{g_A}{2}j_5^Aj_{5A}, \\
\frac{\partial e_H}{\partial\varphi_a^*}&=\big(v_H\varphi\big)_a\qquad(\text{one orbital }\varphi\text{; }\varphi,\varphi^*\text{ independent}), \\
v_H&=\beta\,g_S\rho_S+\beta\,i\chi\,g_P\rho_P+\sum_A\Gamma^0\Gamma^A\,g_Vj_A+\sum_A\Gamma^0\Gamma^A\chi\,g_Aj_{5A}, \\
E&=\sum_{k=1}^{N}\varepsilon_k-\int\Big[(\mathcal S-\mathsf s)\rho_S+(\mathcal P-\mathsf p)\rho_P+(\mathcal V_A-\mathsf v_A)j^A+(\mathcal A_A-\mathsf a_A)j_5^A\Big]d^7x \\
&\quad+E_H+E_{xc}.
\end{align*}
Exchange (Wick’s theorem for contact couplings, $M_{1,2}=\Gamma^0O_{1,2}$) and the Fierz rearrangement:
\begin{align*}
\langle\Phi|:(\bar\psi O_1\psi)(\bar\psi O_2\psi):|\Phi\rangle&=\mathrm{tr}(O_1R)\,\mathrm{tr}(O_2R)-\mathrm{tr}(O_1RO_2R), \\
:\!B(M_1)B(M_2)\!:\;&=B(M_1)B(M_2)-B(M_1M_2), \\
E_x&=-\frac{g_S}2\mathrm{tr}(RR)-\frac{g_P}2\mathrm{tr}(i\chi R\,i\chi R)-\frac{g_V}2\mathrm{tr}(\Gamma^AR\Gamma_AR)-\frac{g_A}2\mathrm{tr}(\Gamma^A\chi R\Gamma_A\chi R), \\
\Sigma_x&=-g_SR-g_P(i\chi)R(i\chi)-g_V\Gamma_AR\Gamma^A-g_A\Gamma_A\chi R\Gamma^A\chi, \\
\delta E_x&=\mathrm{tr}\big(\Sigma_x\,\delta R\big),\qquad E_x=\tfrac12\,\mathrm{tr}\big(\Sigma_xR\big), \\
\mathrm{tr}\big(O_1RO_2R\big)&=\frac1{16}\sum_K\mathrm{tr}\big(\Gamma_K^{-1}R\big)\,\mathrm{tr}\big(O_1\Gamma_KO_2R\big).
\end{align*}
\Needspace{12\baselineskip}\section*{8. The chiral Kohn–Sham theory}\addcontentsline{toc}{section}{8. The chiral Kohn–Sham theory}
{\small
\begin{longtable}{>{\raggedright\arraybackslash}p{\dimexpr 0.1500\linewidth-2\tabcolsep\relax}>{\raggedright\arraybackslash}p{\dimexpr 0.1900\linewidth-2\tabcolsep\relax}>{\raggedright\arraybackslash}p{\dimexpr 0.3600\linewidth-2\tabcolsep\relax}>{\raggedright\arraybackslash}p{\dimexpr 0.3000\linewidth-2\tabcolsep\relax}}
\toprule
\textbf{symbol} & \textbf{Mathematica} & \textbf{definition} & \textbf{meaning} \\
\midrule\endfirsthead
\toprule
\textbf{symbol} & \textbf{Mathematica} & \textbf{definition} & \textbf{meaning} \\
\midrule\endhead
\endfoot
\bottomrule\endlastfoot
$\varphi_{kL}$, $\varphi_{kR}$ & \texttt{orbLV}, \texttt{orbRV} (and \texttt{orbLcV}, \texttt{orbRcV}) & $\varphi_{kL}=E_L^\dagger\varphi_k$, $\varphi_{kR}=E_R^\dagger\varphi_k$ & chiral components of the Kohn–Sham orbitals \\
$\kappa$, $\kappa^*$ & \texttt{κc}, \texttt{κcc} & $\kappa=\sum_k\varphi_{kL}^\dagger\varphi_{kR}$ & complex scalar density \\
$j_L^A$ & \texttt{jLc} & $j_L^A=\sum_k\varphi_{kL}^\dagger\bar\sigma^A\varphi_{kL}$ & left-handed current \\
$j_R^A$ & \texttt{jRc} & $j_R^A=\sum_k\varphi_{kR}^\dagger\sigma^A\varphi_{kR}$ & right-handed current \\
$\varpi$ & — & $\varpi=\mathsf s+i\,\mathsf p$ & complex scalar source \\
$\mathcal M$ & \texttt{Mks} & $\mathcal M=m+\mathcal S+i\,\mathcal P$ & complex Kohn–Sham mass \\
$\mathcal V^L_A$, $\mathcal V^R_A$ & \texttt{VLks}, \texttt{VRks} & $\mathcal V^L_A=\mathcal V_A-\mathcal A_A$, $\mathcal V^R_A=\mathcal V_A+\mathcal A_A$ & chiral vector potentials (lower index) \\
$h_\chi$ & \texttt{hChiKS[p]} & $h_\chi=W^\dagger h_{\rm KS}W$ & chiral Kohn–Sham Hamiltonian, 8 + 8 blocks \\
$\theta$ & — & real angle & chiral ($U(1)_A$) rotation angle \\
$e^{i\theta\chi}$ & \texttt{Uθ} & $e^{i\theta\chi}=\cos\theta\,1_{16}-\sin\theta\,\omega$; \texttt{Uθ} has $\cos\theta=\frac35$, $\sin\theta=\frac45$ & chiral rotation; a REAL matrix \\
$e^{2i\theta}$ & \texttt{e2θ} & $(\frac35+\frac45i)^2$ & phase of $\kappa$ after the rotation \texttt{Uθ} \\*
\end{longtable}}
Equations of the chiral theory:
\begin{align*}
\rho_S&=\kappa+\kappa^*,\qquad \rho_P=i\big(\kappa-\kappa^*\big),\qquad j^A=j_L^A+j_R^A,\qquad j_5^A=j_R^A-j_L^A, \\
\mathsf s\,\rho_S+\mathsf p\,\rho_P&=\varpi\,\kappa+\varpi^*\kappa^*,\qquad \mathsf v_Aj^A+\mathsf a_Aj_5^A=(\mathsf v-\mathsf a)_Aj_L^A+(\mathsf v+\mathsf a)_Aj_R^A, \\
e_H&=\frac{g_S-g_P}{2}\big(\kappa^2+\kappa^{*2}\big)+(g_S+g_P)\,\kappa\kappa^* \\
&\quad+\frac{g_V+g_A}{2}\big(j_L^Aj_{LA}+j_R^Aj_{RA}\big)+(g_V-g_A)\,j_L^Aj_{RA}, \\
h_\chi(\vec p)&=\begin{pmatrix}\mathcal V^L_0-\sigma_i(p^i+\mathcal V^L_i)&\mathcal M\,1_8\\ \mathcal M^*\,1_8&\mathcal V^R_0+\sigma_i(p^i+\mathcal V^R_i)\end{pmatrix}, \\
\mathcal M-m&=\varpi+\frac{\partial(e_H+e_{xc})}{\partial\kappa}, \\
\mathcal V^L_A&=(\mathsf v-\mathsf a)_A+\frac{\partial(e_H+e_{xc})}{\partial j_L^A},\qquad \mathcal V^R_A=(\mathsf v+\mathsf a)_A+\frac{\partial(e_H+e_{xc})}{\partial j_R^A}, \\
e^{i\theta\chi}&:\quad \varphi_L\to e^{-i\theta}\varphi_L,\quad \varphi_R\to e^{i\theta}\varphi_R,\quad \kappa\to e^{2i\theta}\kappa,\quad j_L,\ j_R\ \text{invariant}.
\end{align*}
\Needspace{12\baselineskip}\section*{9. Uniform matter and the Dirac sea}\addcontentsline{toc}{section}{9. Uniform matter and the Dirac sea}
{\small
\begin{longtable}{>{\raggedright\arraybackslash}p{\dimexpr 0.1500\linewidth-2\tabcolsep\relax}>{\raggedright\arraybackslash}p{\dimexpr 0.1900\linewidth-2\tabcolsep\relax}>{\raggedright\arraybackslash}p{\dimexpr 0.3600\linewidth-2\tabcolsep\relax}>{\raggedright\arraybackslash}p{\dimexpr 0.3000\linewidth-2\tabcolsep\relax}}
\toprule
\textbf{symbol} & \textbf{Mathematica} & \textbf{definition} & \textbf{meaning} \\
\midrule\endfirsthead
\toprule
\textbf{symbol} & \textbf{Mathematica} & \textbf{definition} & \textbf{meaning} \\
\midrule\endhead
\endfoot
\bottomrule\endlastfoot
$k_F$ & \texttt{kF}; \texttt{kF0} at saturation & — & Fermi momentum \\
$m^*$ & \texttt{gap7[g, kF]}; \texttt{ms0}, \texttt{msHF} & $m^*=m+\mathcal S$, solution of the gap equation below & effective (Kohn–Sham) mass \\
$E^*$ & — & $E^*=\sqrt{\vec p^{\,2}+m^{*2}}$ & energy of a Kohn–Sham plane wave \\
$|S^6|$ & — & $|S^6|=2\pi^{7/2}/\Gamma_{\rm E}(\tfrac72)=\frac{16\pi^3}{15}$ & area of the unit sphere $S^6\subset\mathbb R^7$ ($\Gamma_{\rm E}$ is Euler’s Gamma function) \\
$c_7$ & \texttt{c7} & $c_7=\frac{8\,|S^6|}{(2\pi)^7}=\frac1{15\pi^4}$ & phase-space factor for 8 spin states in 7 dimensions \\
$\rho$ & \texttt{rho7[kF]}; \texttt{rho0} & $\rho=j^0=\frac{k_F^7}{105\pi^4}$ & particle density \\
$\rho_S$ & \texttt{c7 iS7[ms, kF]}; \texttt{rhoS0} & integral below & scalar density of uniform matter \\
$\varepsilon_{\rm kin}$ & \texttt{c7 iE7[ms, kF]} & integral below & kinetic plus rest energy density \\
$\varepsilon$ & \texttt{eps7[g, gv, kF]} & below & energy density \\
$E/A$ & \texttt{eps7[…]/rho7[kF]} & $E/A=\varepsilon/\rho$ & energy per particle \\
$\varepsilon_F$ & \texttt{fermi7[g, gv, kF]} & $\varepsilon_F=\mathcal V_0+\sqrt{k_F^2+m^{*2}}$ & Fermi energy \\
$\mu$ & \texttt{dEdrho7[g, gv, kF]} & $\mu=d\varepsilon/d\rho$ & chemical potential (numerical derivative with step $10^{-12}$ in the notebook) \\
$\mathcal V_0$ & \texttt{V00} at saturation & $\mathcal V_0=g_V\rho$ (Hartree) & vector potential of uniform matter \\
$G_S$, $G_V$ & \texttt{GSeff}, \texttt{GVeff}; \texttt{GS1}, \texttt{GV1} & below & Fierz-transformed couplings of uniform Hartree–Fock \\
$S_x$, $V_x$ & — & $S_x=\partial E_x/\partial\rho_S$, $V_x=\partial E_x/\partial\rho$ & uniform exchange potentials, $\Sigma_x=S_x1_{16}+V_x\Gamma_0$ \\
$\Lambda$ & \texttt{Λc} & $\Lambda=1$ in the example & momentum cutoff of the Dirac sea, $|\vec p|<\Lambda$ \\
$J(M)$, $K(M)$ & \texttt{Jsea[M]}, \texttt{Ksea[M]} & integrals below & sea integrals \\
$g$ & \texttt{gNJL} & $g=g_S=g_P$; $g=\frac65g_c$ in the example & chirally symmetric coupling \\
$g_c$ & \texttt{gcrit} & $g_c=-\frac1{8J(0)}=-\frac{90\pi^4}{\Lambda^6}$ & critical coupling \\
$h_\chi(\vec p)$ & \texttt{hChiSea} & $\left(\begin{smallmatrix}-\sigma\cdot\vec p&\mathcal M1_8\\\mathcal M^*1_8&\sigma\cdot\vec p\end{smallmatrix}\right)$ & chiral Hamiltonian of the vacuum \\
$\Pi_-(\vec p)$ & \texttt{PiMinus} & $\Pi_-=\frac12\big(1_{16}-h_\chi/E\big)$, $E=\sqrt{\vec p^{\,2}+|\mathcal M|^2}$ & projector on the filled Dirac sea \\
$\mathcal E(\mathcal M)$ & \texttt{Esea[M, m, g]} & below & vacuum energy density \\
$\phi$ & \texttt{phi} (loop variable) & $\mathcal M=|\mathcal M|e^{i\phi}$ & phase of the chiral condensate (Goldstone direction) \\*
\end{longtable}}
Equations of uniform matter (no sea; units $m=1$ in the examples):
\begin{align*}
8\int\frac{d^7k}{(2\pi)^7}\,f(|\vec k|)&=c_7\int_0^\infty k^6f(k)\,dk, \\
\rho_S&=c_7\int_0^{k_F}\frac{k^6\,m^*}{\sqrt{k^2+m^{*2}}}\,dk,\qquad \varepsilon_{\rm kin}=c_7\int_0^{k_F}k^6\sqrt{k^2+m^{*2}}\,dk, \\
m^*&=m+g_S\,\rho_S(m^*),\qquad \varepsilon=\varepsilon_{\rm kin}(m^*)-\frac{(m^*-m)^2}{2g_S}+\frac{g_V}{2}\rho^2,\qquad \mu=\frac{d\varepsilon}{d\rho}=\varepsilon_F, \\
R&=\frac{\rho_S\,1_{16}+\rho\,\Gamma_0}{16}\quad(\text{uniform matter}), \\
G_S&=g_S+\frac{-g_S+g_P-8g_V+8g_A}{16},\qquad G_V=g_V+\frac{-g_S-g_P+6g_V+6g_A}{16} .
\end{align*}
Equations of the Dirac sea (chiral vacuum, $g_P=g_S=g$):
\begin{align*}
J(M)&=\frac{1}{120\pi^4}\int_0^\Lambda\frac{p^6\,dp}{\sqrt{p^2+M^2}},\qquad K(M)=\frac{1}{120\pi^4}\int_0^\Lambda p^6\sqrt{p^2+M^2}\,dp, \\
J(0)&=\frac{\Lambda^6}{720\pi^4},\qquad K(0)=\frac{\Lambda^8}{960\pi^4},\qquad \kappa=-4\,\mathcal M^*J(|\mathcal M|),\qquad \mathcal M=m-8\,g\,J(|\mathcal M|)\,\mathcal M, \\
\mathcal E(\mathcal M)&=-8K(|\mathcal M|)+8J(|\mathcal M|)\,\mathrm{Re}\big[(\mathcal M-m)\mathcal M^*\big]+32\,g\,|\mathcal M|^2J(|\mathcal M|)^2 .
\end{align*}
Numerical parameters of the examples: \texttt{gS0} $=g_S=-10^5$, \texttt{gV0} $=g_V=10^4$ ($g_P=g_A=0$, $m=1$); \texttt{wp} $=50$ (digits of the inner root finding); \texttt{mq0} $=m=\frac1{100}$ (small explicit mass of the sea example); \texttt{M0NJL} and \texttt{MmNJL} are $|\mathcal M|$ for $m=0$ and for $m=\frac1{100}$; \texttt{pS} $=\vec p=(\frac3{10},-\frac15,\frac1{10},\frac14,0,-\frac3{20},\frac1{20})$ and \texttt{Mz7} $=\mathcal M=|\mathcal M|e^{7i/10}$ are the momentum and the complex mass at which $h_\chi(\vec p)$ is diagonalized.

\Needspace{12\baselineskip}\section*{10. Symbols that look alike}\addcontentsline{toc}{section}{10. Symbols that look alike}
{\small
\begin{longtable}{>{\raggedright\arraybackslash}p{\dimexpr 0.3200\linewidth-2\tabcolsep\relax}>{\raggedright\arraybackslash}p{\dimexpr 0.6800\linewidth-2\tabcolsep\relax}}
\toprule
\textbf{symbols} & \textbf{how to tell them apart} \\
\midrule\endfirsthead
\toprule
\textbf{symbols} & \textbf{how to tell them apart} \\
\midrule\endhead
\endfoot
\bottomrule\endlastfoot
$\beta$, $\beta_{\rm b}$, \texttt{βD}, \texttt{β} & $\beta=\Gamma^0$ is the Dirac matrix (\texttt{βD}); $\beta_{\rm b}$ is the rapidity of a boost (the plain symbol \texttt{β} in \texttt{Sβ} and \texttt{Λboost}). \\
$L$, $\mathcal L$, $\Lambda$, \texttt{Λboost}, \texttt{Λc} & $L$ is the 8×8 Lorentz matrix of a boost (Mathematica \texttt{Λboost}); $\mathcal L$ is a Lagrangian density; $\Lambda$ is the momentum cutoff (\texttt{Λc}). The subscript $L$ means left-handed. \\
$\mathcal B(\beta_{\rm b})$, $B(M)$ & $\mathcal B$ is the spinor boost; $B(M)$ the one-body operator $\sum M_{ab}c_a^\dagger c_b$; $B$ is also a spacetime index. \\
$\mathcal S$, $\mathsf s$, $s$, $S_x$, $S^6$ & $\mathcal S$ scalar Kohn–Sham potential; $\mathsf s$ scalar source; $s$ spin label; $S_x$ uniform scalar exchange potential; $S^6$ unit sphere. \\
$\mathsf p$, $p$, $\vec p$ & $\mathsf p$ pseudoscalar source; $p=(E,\vec p)$ 8-momentum; $\vec p$ momentum. \\
$\mathsf a_A$, $\mathcal A_A$, $a$ & axial source; axial Kohn–Sham potential; spinor index. \\
$\mathsf v_A$, $\mathcal V_A$, $V$, $V_x$ & vector source; vector Kohn–Sham potential; box volume; uniform vector exchange potential. \\
$\mu$, $\varpi$ & $\mu$ is the chemical potential $d\varepsilon/d\rho$; $\varpi=\mathsf s+i\mathsf p$ is the complex scalar source. \\
$M(\vec p)$, $\mathcal M$, $M$ in $J(M)$, $M_{1,2}$ & real 8×8 block of $h(\vec p)$; complex Kohn–Sham mass; argument of the sea integrals; matrices $\Gamma^0O_{1,2}$ of Wick’s theorem. \\
$K$ in $\Gamma_K$, $K(M)$, index $K$ & ordered index set of a basis element; sea integral; spacetime index. \\
$J(M)$, $j^A$, $j_5^A$, $j_L^A$, $j_R^A$ & sea integral; vector, axial, left and right currents. \\
$\kappa$ (field), $\kappa$ (Kohn–Sham) & the same bilinear $\phi_L^\dagger\phi_R$, evaluated for the field or summed over the Kohn–Sham orbitals. \\
$\phi_L,\phi_R$, $\varphi_k$, $\varphi_{kL},\varphi_{kR}$, $\phi$ & chiral fields; Kohn–Sham orbitals; their chiral components; phase of $\mathcal M$. \\
$\theta_k$, $\theta$ & Grassmann generators; chiral rotation angle. \\
$\sigma$, $\sigma_i$, $\sigma^A$, $\bar\sigma^A$ & 8×8 matrix $\tau_1\tau_2\tau_3$; Pauli matrices $i\tau^{O}_i$; chiral 8×8 matrices. \\
$\Sigma_{AB}$, $\Sigma_x$, $\Sigma_H$ & Lorentz generators; exchange self-energy; Hartree self-energy. \\
$\Gamma_A$, $\Gamma_{\rm E}$ & Dirac matrices; Euler’s Gamma function (only in $|S^6|$). \\
$E$, $E^*$, $E_L$, $E_R$, $E_H$, $E_x$, $E_{xc}$, $E[w]$, $E/A$, $\mathcal E$ & free energy $\sqrt{m^2+\vec p^{\,2}}$; Kohn–Sham plane-wave energy; chiral isometries; Hartree, exchange, exchange-correlation energies; ground-state energy; energy per particle; vacuum energy density. \\
$\varepsilon$, $\varepsilon_k$, $\varepsilon_F$, $\varepsilon_{\rm kin}$ & energy density; Kohn–Sham eigenvalues; Fermi energy; kinetic energy density. \\
$e_s$, $e_H$, $e_i$ & unit column of length 8; Hartree energy density; octonion units. \\
$\rho$, $\rho_S$, $\rho_P$, $\varrho$ & particle density $j^0$; scalar and pseudoscalar densities; one-body density matrix. \\
$R$, $\varrho$ & $R=\varrho\,\Gamma^0$, both 16×16. \\
$h$, $h(\vec p)$, $h_{\rm KS}$, $h_\chi$ & Dirac Hamiltonian; its plane-wave form; Kohn–Sham Hamiltonian; its chiral form. \\
$H$, $\mathcal H$, $H[w]$, $\mathbb H$ & Hamiltonian; Hamiltonian density; Hamiltonian with sources; quaternions. \\
$Q$, $Q(v)$ & charge; quadratic form $v\cdot v$ of the Clifford algebra. \\
$T$, $T_A$, $T_s$ & invertible matrix of the proof in section 4; Dirac matrices of 4+4.nb; Kohn–Sham kinetic energy. \\
$Y$, $Z$, $R_A$, $R$ & matrices of the two proofs in section 4; the Kohn–Sham density matrix $R$. \\
$\Gamma_A$ and $T_A$ & $\Gamma_0=T_0$, $\Gamma_i=T_0T_i$ for $i=1,2,3$, $\Gamma_k=T_0T_kT_4T_5T_6T_7$ for $k=4,\dots,7$. \\
$c_7$, $c_a$ & phase-space factor; mode operators. \\
$w$, $W$ & collection of sources; chiral basis change. \\
$N$, $n$ & number of particles; collection of densities. \\
$\delta_{ij}$, $\delta^7$, $\delta E/\delta\rho$ & Kronecker delta; Dirac delta; functional derivative. \\
$g$, $g_S,g_P,g_V,g_A$, $g_c$, $G_S,G_V$ & common coupling $g_S=g_P$ of the sea example; the four couplings; critical coupling; Fierz-transformed couplings. \\*
\end{longtable}}
\Needspace{12\baselineskip}\section*{11. Every Mathematica name of sections 7–10 of 7+1.nb}\addcontentsline{toc}{section}{11. Every Mathematica name of sections 7–10 of 7+1.nb}
The names are listed in the order of the notebook sections in which they are defined. Names defined in sections 0–6 of the notebook and used here: \texttt{X}, \texttt{η17}, \texttt{η17inv}, \texttt{ID8}, \texttt{ID16}, \texttt{Zero8}, \texttt{Zero16}, \texttt{zero8}, \texttt{τ[A]}, \texttt{τO[i]}, \texttt{σ}, \texttt{Γ17[A]} ($A=0,\dots,8$; \texttt{Γ17[8]} $=\omega$), \texttt{Γ17up[A]}, \texttt{products256}, \texttt{χ17}, \texttt{PL17}, \texttt{PR17}, \texttt{Σ17[A, B]}, \texttt{Sβ}, \texttt{Sβinv}, \texttt{Λboost}, \texttt{DiracOp17}, \texttt{waveOp17}, \texttt{check}, \texttt{checks}. The symbols \texttt{mD} ($m$), \texttt{p1}, …, \texttt{p7} ($p^i$), \texttt{gS}, \texttt{gP}, \texttt{gV}, \texttt{gA}, \texttt{Sks}, \texttt{Pks}, \texttt{Vks[A]}, \texttt{Aks[A]}, \texttt{Vbox}, \texttt{β}, \texttt{ρSsym}, \texttt{ρsym}, \texttt{V0s}, \texttt{mst}, \texttt{εs} are left without values (symbolic).

\Needspace{8\baselineskip}\subsection*{11.1 Section 7: the Dirac equation of dirac16complex}\addcontentsline{toc}{subsection}{11.1 Section 7: the Dirac equation of dirac16complex}
{\footnotesize
\begin{longtable}{>{\raggedright\arraybackslash}p{\dimexpr 0.2100\linewidth-2\tabcolsep\relax}>{\raggedright\arraybackslash}p{\dimexpr 0.4700\linewidth-2\tabcolsep\relax}>{\raggedright\arraybackslash}p{\dimexpr 0.3200\linewidth-2\tabcolsep\relax}}
\toprule
\textbf{name} & \textbf{Wolfram Language definition (abbreviated)} & \textbf{meaning} \\
\midrule\endfirsthead
\toprule
\textbf{name} & \textbf{Wolfram Language definition (abbreviated)} & \textbf{meaning} \\
\midrule\endhead
\endfoot
\bottomrule\endlastfoot
\texttt{βD}, \texttt{αD[i]}, \texttt{ω17} & \texttt{βD = Γ17up[0]}; \texttt{αD[i] = Γ17up[0].Γ17up[i]}; \texttt{ω17 = Γ17[8]} & $\beta$, $\alpha_i$, $\omega$ \\
\texttt{dirac16complex} & \texttt{Table[ψ16[a] @@ X, \{a, 1, 16\}]} & $\psi$: 16 arbitrary functions of $x_0,\dots,x_7$ \\
\texttt{DiracEq16[ψ, m]} & \texttt{I DiracOp17[ψ] - m ψ} & $(i\Gamma^A\partial_A-m)\psi$ \\
\texttt{hD[ψ, m]} & \texttt{Sum[-I αD[i].D[ψ, X[[i+1]]], \{i,1,7\}] + m βD.ψ} & $h\psi$ \\
\texttt{pv7}, \texttt{asm7} & \texttt{\{p1, …, p7\}}; \texttt{mD > 0 \&\& p1, …, p7 real} & $\vec p$; assumptions for \texttt{Simplify} \\
\texttt{Ep}, \texttt{Mp} & \texttt{Sqrt[mD\textasciicircum{}2 + pv7.pv7]}; \texttt{mD ID8 + Sum[pv7[[i]] τO[i], \{i,1,7\}]} & $E$, $M(\vec p)$ \\
\texttt{hp[q, m]} & \texttt{Sum[q[[i]] αD[i], \{i,1,7\}] + m βD} & $h(\vec q)$ \\
\texttt{uS[s, M, E]}, \texttt{vS[s, M, E]} & \texttt{Join[E e\_s, Transpose[M].e\_s]/Sqrt[E]}; \texttt{Join[E e\_s, -M.e\_s]/Sqrt[E]} & $u_s$, $v_s$ \\
\texttt{slash17[E, q]} & \texttt{E Γ17[0] + Sum[q[[i]] Γ17[i], \{i,1,7\}]} & $\not p$ \\
\texttt{p0v}, \texttt{M0}, \texttt{u0Matrix} & \texttt{\{2, 2, 4, 0, 0, 0, 0\}}; \texttt{ID8 + p0v.τO}; \texttt{Transpose[Table[uS[s, M0, 5], \{s,1,8\}]]} & example momentum ($m=1$, $E=5$), $M$, the matrix $(u_1,\dots,u_8)$ \\
\texttt{USmat}, \texttt{VSmat}, \texttt{VSmatMinus} & columns $u_s(\vec p)$, $v_s(\vec p)$, $v_s(-\vec p)$ (with $M(-\vec p)=M^{\mathsf T}$) & 16×8 matrices of spinors \\
\texttt{Umat} & \texttt{Join[USmat, VSmatMinus, 2]/Sqrt[2 Ep]} & $U(\vec p)$ \\
\texttt{grass[\{k1, …\}]} & ordered monomial (inert head) & $\theta_{k_1}\cdots\theta_{k_r}$ \\
\texttt{gTerms}, \texttt{gSplit} & list of the terms of an expanded sum; \{coefficient, monomial\} of a term & helpers of \texttt{gMul}, \texttt{gD} \\
\texttt{gProd} & \texttt{0} if two monomials share a generator, else \texttt{Signature[...] grass[Sort[...]]} & product of monomials \\
\texttt{gMul[x, y]} & bilinear extension of \texttt{gProd} & Grassmann product \\
\texttt{gD[k, x]} & left derivative & $\partial^{L}x/\partial\theta_k$ \\
\texttt{θg[k]} & \texttt{grass[\{k\}]} & $\theta_k$ \\
\texttt{conjR[e]} & \texttt{e /. Complex[r, s] :> Complex[r, -s]} & complex conjugate of an expression whose symbols are real \\
\texttt{mom5}, \texttt{E5}, \texttt{M5} & momenta $\pm(2,2,4,0,0,0,0)$, $\pm(0,0,0,2,2,0,0)$, $\vec0$; energies $5,5,3,3,1$; $M(\vec p_n)$ & five modes ($m=1$) \\
\texttt{Ugr} & \texttt{\{uS[1, M5[[1]], 5] e\textasciicircum{}\{-i(5 x0 - p1.x)\}, vS[3, M5[[3]], 3] e\textasciicircum{}\{i(3 x0 - p3.x)\}, uS[5, M5[[5]], 1] e\textasciicircum{}\{-i x0\}\}} & $U_1,U_2,U_3$ \\
\texttt{ψG}, \texttt{ψbarG}, \texttt{jG} & $\sum_k\theta_kU_k$; $\sum_k\theta_{3+k}U_k^\dagger\Gamma^0$; $\bar\psi_G\Gamma^A\psi_G$ with \texttt{gMul} & Grassmann-valued solution, adjoint, current \\
\texttt{xv7}, \texttt{yv7} & \texttt{Rest[X]}; \texttt{Array[y7, 7]} & $\vec x$, $\vec y$ \\
\texttt{modeB[n, s, x]}, \texttt{modeD[n, s, x]} & \texttt{uS[…] Exp[-I E x0 + I p.x]/Sqrt[2 E Vbox]}; \texttt{vS[…] Exp[I E x0 - I p.x]/Sqrt[2 E Vbox]} & coefficients of $b_s(\vec p_n)$ and $d_s^\dagger(\vec p_n)$ in $\psi$ \\
\texttt{carMatrix} & $\sum_{n,s}$ [\texttt{modeB} \texorpdfstring{\ensuremath{\otimes}}{⊗} \texttt{conjR[modeB]} + \texttt{modeD} \texorpdfstring{\ensuremath{\otimes}}{⊗} \texttt{conjR[modeD]}] & $\lbrace\psi_a(\vec x),\psi_b^\dagger(\vec y)\rbrace$ from the mode algebra \\
\texttt{ψreF}, \texttt{ψimF} & 16 arbitrary real functions each & real and imaginary parts of $\psi$ (charge-conjugation check) \\
\texttt{Sb0}, \texttt{Λb0}, \texttt{pPrime} & \texttt{Sβ}, \texttt{Λboost} at $\beta_{\rm b}=2\log2$; \texttt{Λb0.\{5, 2, 2, 4, 0, 0, 0, 0\}} & $\mathcal B$, $L$, $p'=(\frac{55}8,-\frac{41}8,2,4,0,0,0,0)$ \\
\texttt{EL}, \texttt{ER}, \texttt{Wch} & \texttt{Join[ID8, -I ID8]/Sqrt[2]}; \texttt{Join[-I ID8, ID8]/Sqrt[2]}; \texttt{Join[EL, ER, 2]} & $E_L$, $E_R$, $W$ \\
\texttt{ϕL}, \texttt{ϕR} & \texttt{ConjugateTranspose[EL].dirac16complex}; same with \texttt{ER} & $\phi_L$, $\phi_R$ \\
\texttt{σW[i]}, \texttt{σWp[q]} & \texttt{I τO[i]}; \texttt{Sum[q[[i]] σW[i], \{i,1,7\}]} & $\sigma_i$, $\sigma\cdot\vec q$ \\
\texttt{tauD[f]} & \texttt{Sum[τO[i].D[f, X[[i+1]]], \{i,1,7\}]} & $\tau^{O}_i\partial_if$ \\
\texttt{f8} & \texttt{Table[f8c[a] @@ X, \{a, 1, 8\}]} & 8 arbitrary functions (Klein–Gordon check of the chiral equations) \\*
\end{longtable}}
\Needspace{8\baselineskip}\subsection*{11.2 Section 8: the 16-component Kohn–Sham theory}\addcontentsline{toc}{subsection}{11.2 Section 8: the 16-component Kohn–Sham theory}
{\footnotesize
\begin{longtable}{>{\raggedright\arraybackslash}p{\dimexpr 0.2100\linewidth-2\tabcolsep\relax}>{\raggedright\arraybackslash}p{\dimexpr 0.4700\linewidth-2\tabcolsep\relax}>{\raggedright\arraybackslash}p{\dimexpr 0.3200\linewidth-2\tabcolsep\relax}}
\toprule
\textbf{name} & \textbf{Wolfram Language definition (abbreviated)} & \textbf{meaning} \\
\midrule\endfirsthead
\toprule
\textbf{name} & \textbf{Wolfram Language definition (abbreviated)} & \textbf{meaning} \\
\midrule\endhead
\endfoot
\bottomrule\endlastfoot
\texttt{vPairs} & \texttt{\{\{gS, ID16, ID16\}, \{gP, I χ17, I χ17\}\}} + \texttt{\{gV, Γ17up[A], Γ17[A]\}} + \texttt{\{gA, Γ17up[A].χ17, Γ17[A].χ17\}} & the 18 vertex triples $(g,O,O')$ \\
\texttt{ψP}, \texttt{ψbarP} & \texttt{grass[\{a\}]}; \texttt{grass[\{16 + a\}]} & 32 Grassmann generators \\
\texttt{gBil[M]} & $\sum_{a,b}M_{ab}$ \texttt{gMul[ψbarP[[a]], ψP[[b]]]} & $\bar\psi M\psi$ \\
\texttt{LintG} & \texttt{Sum[-g/2 gMul[gBil[O], gBil[Oc]], \{g, O, Oc\} in vPairs]} & $\mathcal L_{\rm int}$ \\
\texttt{VKS}, \texttt{AKS} & \texttt{Array[Vks, 8, 0]}; \texttt{Array[Aks, 8, 0]} & $\mathcal V_A$, $\mathcal A_A$ (lower index) \\
\texttt{hKS[q]} & \texttt{Sum[αD[i] (q[[i]] + VKS[[i+1]])] + βD (mD + Sks) + VKS[[1]] ID16 + βD.(I χ17) Pks + χ17 AKS[[1]] + Sum[αD[i].χ17 AKS[[i+1]]]} & $h_{\rm KS}(\vec q)$ \\
\texttt{orbV}, \texttt{orbcV} & \texttt{Array[orb, 16]}; \texttt{Array[orbc, 16]} & $\varphi$ and $\varphi^*$ (independent variables) \\
\texttt{dens16[M]} & \texttt{orbcV.M.orbV} & $\varphi^\dagger M\varphi$ \\
\texttt{EHfun[d]}, \texttt{EH16} & $\frac{g_S}2d[\Gamma^0]^2+\frac{g_P}2d[\Gamma^0i\chi]^2+\frac{g_V}2\sum_A\eta_{AA}d[\Gamma^0\Gamma^A]^2+\frac{g_A}2\sum_A\eta_{AA}d[\Gamma^0\Gamma^A\chi]^2$; \texttt{EHfun[dens16]} & $e_H$ for a density function \texttt{d}; for one orbital \\
\texttt{vH16} & $\beta g_S\rho_S+\beta i\chi g_P\rho_P+\sum_A\eta_{AA}\Gamma^0\Gamma^Ag_Vj^A+\sum_A\eta_{AA}\Gamma^0\Gamma^A\chi g_Aj_5^A$ & $v_H$ \\
\texttt{basis3}, \texttt{idx3} & \texttt{Subsets[Range[16], \{3\}]}; position of each subset & the 560 three-particle basis states \\
\texttt{cdc[a, b]} & sparse 560×560 matrix with signs from the anticommutation & $c_a^\dagger c_b$ on three-particle states \\
\texttt{Bop[M]} & \texttt{Sum[M[[a, b]] cdc[a, b]]} & $B(M)$ \\
\texttt{gsOrth} & Gram–Schmidt orthonormalization of a list of vectors & helper \\
\texttt{orbs3}, \texttt{Φ3} & 3 random orthonormal orbitals (\texttt{SeedRandom[71]}, 40 digits); their Slater determinant & state of the Wick check \\
\texttt{Rorb} & \texttt{Transpose[orbs3].Conjugate[orbs3].Γ17[0]} & $R=\sum_k\varphi_k\bar\varphi_k$ \\
\texttt{wickLHS[o1, o2]} & $\langle\Phi|B(M_1)B(M_2)|\Phi\rangle-\langle\Phi|B(M_1M_2)|\Phi\rangle$, $M_i=\Gamma^0o_i$ & left side of Wick’s theorem \\
\texttt{Rr}, \texttt{Dr} & random complex integer 16×16 matrices (entries $-3,\dots,3$) & test matrices for $R$ and $\delta R$ \\
\texttt{Exfun[R]}, \texttt{Sigmax[R]} & \texttt{Sum[-g/2 Tr[O.R.Oc.R]]}; \texttt{Sum[-g Oc.R.O]} over \texttt{vPairs} & $E_x$, $\Sigma_x$ \\
\texttt{F256} & \texttt{Flatten /@ products256} & the 256 basis matrices as vectors (Fierz completeness) \\
\texttt{Runi}, \texttt{ExUni} & \texttt{(ρSsym ID16 + ρsym Γ17[0])/16}; the formula of $E_x$ for uniform matter & $R$ and $E_x$ of uniform matter \\
\texttt{GSeff[gs, gp, gv, ga]}, \texttt{GVeff[…]} & \texttt{gs + (-gs + gp - 8 gv + 8 ga)/16}; \texttt{gv + (-gs - gp + 6 gv + 6 ga)/16} & $G_S$, $G_V$ \\
\texttt{c7}, \texttt{rho7[kF]} & \texttt{1/(15 Pi\textasciicircum{}4)}; \texttt{kF\textasciicircum{}7/(105 Pi\textasciicircum{}4)} & $c_7$, $\rho$ \\
\texttt{iS7[ms, kF]}, \texttt{iE7[ms, kF]} & \texttt{Integrate[k\textasciicircum{}6 ms/Sqrt[k\textasciicircum{}2 + ms\textasciicircum{}2], \{k, 0, kF\}]}; \texttt{Integrate[k\textasciicircum{}6 Sqrt[k\textasciicircum{}2 + ms\textasciicircum{}2], \{k, 0, kF\}]} (closed forms) & $\rho_S/c_7$, $\varepsilon_{\rm kin}/c_7$ \\
\texttt{gS0}, \texttt{gV0}, \texttt{wp} & \texttt{-10\textasciicircum{}5}, \texttt{10\textasciicircum{}4}, \texttt{50} & couplings of the example; working precision \\
\texttt{gap7[g, kF]} & root of \texttt{ms - 1 - g c7 iS7[ms, kF]} (Brent, 50 digits) & $m^*$ \\
\texttt{eps7[g, gv, kF]} & \texttt{c7 iE7[m*, kF] - (m* - 1)\textasciicircum{}2/(2 g) + gv/2 rho7[kF]\textasciicircum{}2} & $\varepsilon$ \\
\texttt{fermi7[g, gv, kF]} & \texttt{gv rho7[kF] + Sqrt[kF\textasciicircum{}2 + m*\textasciicircum{}2]} & $\varepsilon_F$ \\
\texttt{dEdrho7[g, gv, k0]} & symmetric difference quotient with $h=10^{-12}$ & $d\varepsilon/d\rho$ \\
\texttt{kF0} & root of \texttt{eps7/rho7 - fermi7} in $[0.7,0.8]$ (40 digits) & saturation $k_F$ \\
\texttt{ms0}, \texttt{rho0}, \texttt{rhoS0}, \texttt{V00}, \texttt{eps0} & \texttt{gap7[gS0, kF0]}, \texttt{rho7[kF0]}, \texttt{c7 iS7[ms0, kF0]}, \texttt{gV0 rho0}, \texttt{eps7[gS0, gV0, kF0]} & saturation values \\
\texttt{hartreeSaturation}, \texttt{hartreeKF1} & rules with the printed numbers & tables of section 8.7 of the notebook \\
\texttt{GS1}, \texttt{GV1}, \texttt{msHF}, \texttt{rhoHF}, \texttt{rhoSHF} & \texttt{GSeff[gS0, 0, gV0, 0]}, \texttt{GVeff[…]}; $m^*$, $\rho$, $\rho_S$ at $k_F=\frac34$ & Hartree–Fock example \\
\texttt{hartreeFock34}, \texttt{eAcurve[g, gv]} & rules with the printed numbers; list of $(k_F, E/A-1)$ & Hartree–Fock table; data of the plot \\*
\end{longtable}}
\Needspace{8\baselineskip}\subsection*{11.3 Section 9: the chiral Kohn–Sham theory}\addcontentsline{toc}{subsection}{11.3 Section 9: the chiral Kohn–Sham theory}
{\footnotesize
\begin{longtable}{>{\raggedright\arraybackslash}p{\dimexpr 0.2100\linewidth-2\tabcolsep\relax}>{\raggedright\arraybackslash}p{\dimexpr 0.4700\linewidth-2\tabcolsep\relax}>{\raggedright\arraybackslash}p{\dimexpr 0.3200\linewidth-2\tabcolsep\relax}}
\toprule
\textbf{name} & \textbf{Wolfram Language definition (abbreviated)} & \textbf{meaning} \\
\midrule\endfirsthead
\toprule
\textbf{name} & \textbf{Wolfram Language definition (abbreviated)} & \textbf{meaning} \\
\midrule\endhead
\endfoot
\bottomrule\endlastfoot
\texttt{orbLV}, \texttt{orbLcV}, \texttt{orbRV}, \texttt{orbRcV} & \texttt{Array[orbL, 8]}, \texttt{Array[orbLc, 8]}, \texttt{Array[orbR, 8]}, \texttt{Array[orbRc, 8]} & $\varphi_L$, $\varphi_L^*$, $\varphi_R$, $\varphi_R^*$ \\
\texttt{psiChi}, \texttt{psiChiDag} & \texttt{EL.orbLV + ER.orbRV}; \texttt{orbLcV.EL\textasciicircum{}† + orbRcV.ER\textasciicircum{}†} & $\varphi=E_L\varphi_L+E_R\varphi_R$ and $\varphi^\dagger$ \\
\texttt{densChi[M]} & \texttt{psiChiDag.M.psiChi} & $\varphi^\dagger M\varphi$ in chiral variables \\
\texttt{κc}, \texttt{κcc} & \texttt{orbLcV.orbRV}; \texttt{orbRcV.orbLV} & $\kappa$, $\kappa^*$ \\
\texttt{σUp}, \texttt{σbarUp} & \texttt{\{ID8, σW[1], …, σW[7]\}}; \texttt{\{ID8, -σW[1], …, -σW[7]\}} & $\sigma^A$, $\bar\sigma^A$ \\
\texttt{jLc}, \texttt{jRc} & \texttt{orbLcV.σbarUp[[A+1]].orbLV}; \texttt{orbRcV.σUp[[A+1]].orbRV} & $j_L^A$, $j_R^A$ \\
\texttt{EHchiAbs[k, kc, jl, jr]}, \texttt{EHchi} & $\frac{g_S}2(k+k_c)^2-\frac{g_P}2(k-k_c)^2+\frac{g_V+g_A}2(j_l\cdot j_l+j_r\cdot j_r)+(g_V-g_A)\,j_l\cdot j_r$; at \texttt{κc}, \texttt{κcc}, \texttt{jLc}, \texttt{jRc} & $e_H$ in chiral variables \\
\texttt{VLks}, \texttt{VRks}, \texttt{Mks} & \texttt{VKS - AKS}, \texttt{VKS + AKS}, \texttt{mD + Sks + I Pks} & $\mathcal V^L_A$, $\mathcal V^R_A$, $\mathcal M$ \\
\texttt{hChiKS[q]} & \texttt{\{\{VL0 ID8 - Sum[σW[i] (q[[i]] + VLi)], Mks ID8\}, \{(mD + Sks - I Pks) ID8, VR0 ID8 + Sum[σW[i] (q[[i]] + VRi)]\}\}} & $h_\chi(\vec q)$ \\
\texttt{VLc}, \texttt{VRc} & $(g_V+g_A)j_{LA}+(g_V-g_A)j_{RA}$; $(g_V+g_A)j_{RA}+(g_V-g_A)j_{LA}$ & Hartree $\mathcal V^L_A$, $\mathcal V^R_A$ \\
\texttt{dEHdk}, \texttt{dEHdkc} & $g_S(\kappa+\kappa^*)-g_P(\kappa-\kappa^*)$; $g_S(\kappa+\kappa^*)+g_P(\kappa-\kappa^*)$ & $\partial e_H/\partial\kappa$, $\partial e_H/\partial\kappa^*$ \\
\texttt{Uθ}, \texttt{e2θ} & \texttt{3/5 ID16 + I 4/5 χ17}; \texttt{(3/5 + 4/5 I)\textasciicircum{}2} & $e^{i\theta\chi}$ with $\cos\theta=\frac35$; $e^{2i\theta}$ \\
\texttt{EHrot} & $e_H$ of the rotated orbital $e^{i\theta\chi}\varphi$ & chiral-rotation check \\
\texttt{κ16}, \texttt{κ16c} & \texttt{dens16[EL.ER\textasciicircum{}†]}; \texttt{dens16[ER.EL\textasciicircum{}†]} & $\kappa$, $\kappa^*$ from a 16-component orbital \\
\texttt{kappaUni}, \texttt{jLUni}, \texttt{jRUni} & traces with \texttt{Runi} & $\kappa$, $j_L^A$, $j_R^A$ of uniform matter \\
\texttt{hChiUni} & \texttt{\{\{V0s ID8 - σWp[pv7], mst ID8\}, \{mst ID8, V0s ID8 + σWp[pv7]\}\}} & $h_\chi$ of uniform matter ($\mathcal V_0$, $m^*$ symbolic) \\
\texttt{chiralUniform} & rules with the printed numbers & table of section 9.5 of the notebook \\
\texttt{Λc}, \texttt{Jsea[M]}, \texttt{Ksea[M]} & \texttt{1}; closed forms of $J(M)$, $K(M)$ & cutoff and sea integrals \\
\texttt{J0sea}, \texttt{K0sea}, \texttt{gcrit} & $J(0)$, $K(0)$; \texttt{-1/(8 J0sea)} & $J(0)$, $K(0)$, $g_c$ \\
\texttt{Esea[Mz, mq, g]}, \texttt{E0sea} & \texttt{-8 Ksea[|Mz|] + 8 Jsea[|Mz|] Re[(Mz - mq) Mz\textasciicircum{}*] + 32 g |Mz|\textasciicircum{}2 Jsea[|Mz|]\textasciicircum{}2}; \texttt{-8 K0sea} & $\mathcal E(\mathcal M)$; $\mathcal E(0)$ for $m=0$ \\
\texttt{gNJL}, \texttt{gapNJL[M, mq]} & \texttt{6/5 gcrit}; \texttt{M - mq + 8 gNJL Jsea[M] M} & $g$; the gap equation for real $\mathcal M>0$ \\
\texttt{M0NJL}, \texttt{mq0}, \texttt{MmNJL} & root of \texttt{gapNJL[M, 0]/M}; \texttt{1/100}; root of \texttt{gapNJL[M, mq0]} & $|\mathcal M|$ for $m=0$; $m$; $\mathcal M$ for $m=\frac1{100}$ \\
\texttt{njlVacuum} & rules with the printed numbers & table of section 9.6 of the notebook \\
\texttt{pS}, \texttt{Mz7}, \texttt{Esp} & $(\frac3{10},-\frac15,\frac1{10},\frac14,0,-\frac3{20},\frac1{20})$; \texttt{M0NJL Exp[7 I/10]}; $\sqrt{\vec p^{\,2}+|\mathcal M|^2}$ & test momentum, complex mass, energy \\
\texttt{hChiSea}, \texttt{PiMinus}, \texttt{h0Chi} & $h_\chi(\vec p)$; $\frac12(1_{16}-h_\chi/E)$; $h_\chi$ with $\mathcal M=0$ & sea Hamiltonian, projector, massless Hamiltonian \\*
\end{longtable}}
Section 10 of the notebook defines no new names; it prints \texttt{Length[checks]}, \texttt{failedChecks}, \texttt{allChecksPassed} and a table of all checks.

\end{document}
